\chapter{Evaluationsdurchführung}
\label{chap:evaluationsdurchfuehrung}

\section{Durchführung der Tests}
\label{sec:durchfuehrung_tests}

Zur Durchführung der Tests wird für jedes untersuchte Modell in Microsoft Copilot Studio ein eigener Agent angelegt. Alle Agenten erhielten dieselben Anweisungen, das selbe hinterlegte Wissen und die selben Konfigurationen, sodass sich die Agenten ausschließlich im zugrunde liegenden Sprachmodell unterscheiden. Der Prompt wurde zuvor anhand von fünf separaten Rechnungen mit Freigaben entwickelt (vgl. Abschnitt~\ref{sec:prompt}) und vor Beginn der Tests eingefroren. Damit wird vermieden, dass der Prompt auf die Testdaten optimiert wird und die Ergebnisse dadurch zu positiv ausfallen. Der finale Prompt befindet sich in Anhang \ref{anh:prompt}.

Untersucht wurden die Modelle Opus 4.8, Opus 4.7, Opus 4.6, Sonnet 4.6, GPT 5.5, GPT 5.3, GPT 5 und GPT 4.1. Das sind alle Modelle, die zu diesem Zeitpunkt (20.10.2026) im Copilot Studio verfügbar sind.

Die Durchführung erfolgte automatisiert über einen Flow in Power Automate. Jeder Agent erhält die 100 Rechnungen selbst und die zugehörige Freigabe als Anhang übergeben. Dann wird der Agent 100 mal ausgeführt und die Antworten zusammen mit Rechnungsnummer, Zeitstempel, Modell und Laufnummer in eine Excel Tabelle geschrieben. 

Die Auswertung erfolgt anschließend in Excel. Dort wird für jede Antwort festgehalten, ob das Modell eine Exception gemeldet hat und welchen Grund es nennt. Diese Angaben werden über die Rechnungsnummer mit der Abgleichstabelle (Ground Truth) verknüpft. Aus dem Abgleich ergibt sich für jede Rechnung, ob es sich um ein \ac{TP}, \ac{FP}, \ac{TN} oder \ac{FN} handelt. Abschließend werden diese Werte je Modell gezählt und bilden die Grundlage für die Berechnung der in Abschnitt~\ref{sec:bewertungsmatrix} definierten Kennzahlen.

Für die Einordnung als \ac{TP} reicht es nicht aus, dass eine Rechnung als Exception markiert wurde. Zusätzlich wird der vom Modell genannte Grund für die Exception mit dem in der Ground Truth hinterlegten Grund verglichen. Nur wenn beide übereinstimmen, wird die Antwort als \ac{TP} gewertet. Weicht der Grund ab, so wird die Antwort als \ac{FN} gewertet, da der eingebaute Fehler in diesem Fall unentdeckt bliebe. Auch die Antworten für \ac{TN}, \ac{FP}, \ac{FN} werden überprüft, um technische Fehler oder anderweiterige Probleme auszuschließen.

Zur Messung der Konsistenz wurden 30 der 100 Rechnungen von jedem Modell weitere viermal ausgewertet, sodass für diese Rechnungen insgesamt fünf Durchläufe je Modell vorliegen. Die Auswahl der Rechnungen ist in Abschnitt~\ref{sec:bewertungsmatrix} beschrieben. Die Konsistenz ergibt sich aus dem Mittelwert der Anteile, mit denen die einzelnen Befunde über die fünf Durchläufe hinweg gemeldet wurden. Precision und Recall beruhen dagegen ausschließlich auf dem ersten Durchlauf über alle 100 Rechnungen.

Während der Durchführung zeigten sich zwei technische Auffälligkeiten. Erstens wurden die Freigabescreenshots in den Rechnungsunterlagen häufig nicht erkannt und konnten deshalb nicht ausgelesen werden. Zweitens werden PNG Dateien im Chatfenster von Copilot Studio gar nicht verarbeitet. Somit ist eine Nutzung  Diese Einschränkungen betreffen alle Modelle gleichermaßen und werden bei der Interpretation der Ergebnisse berücksichtigt.

\section{Ergebnisse der Stichprobenanalysen}
\label{sec:ergebnisse_stichproben}


Der Testdatensatz umfasst 100 Rechnungen mit den zugehörigen Freigaben. Laut Ground Truth enthalten 24 davon eine Exception.

\subsection{Ergebnisse der Modelle}
\label{subsec:ergebnisse_modelle}

Tabelle~\ref{tab:konfusion_rechnungen} zeigt die Verteilung der Ergebnisse je Modell, Tabelle~\ref{tab:metriken_rechnungen} die daraus berechneten Kennzahlen.

\begin{table}[htbp]
  \centering
  \caption{Konfusionswerte je Modell}
  \label{tab:konfusion_rechnungen}
  \begin{tabular}{lcccc}
    \hline
    Modell & TP & FP & TN & FN \\
    \hline
    Opus 4.7   & 22 &  1 & 75 &  2 \\
    Opus 4.8   & 20 &  2 & 74 &  4 \\
    Opus 4.6   & 21 & 12 & 64 &  3 \\
    Sonnet 4.6 & 20 & 12 & 64 &  4 \\
    GPT 4.1    & 13 &  2 & 74 & 11 \\
    GPT 5.5    & 19 & 21 & 55 &  5 \\
    GPT 5.3    & 18 & 21 & 55 &  6 \\
    GPT 5      & 13 & 16 & 60 & 11 \\
    \hline
  \end{tabular}
\end{table}

\begin{table}[htbp]
  \centering
  \caption{Kennzahlen je Modell (in \%)}
  \label{tab:metriken_rechnungen}
  \begin{tabular}{lccccc}
    \hline
    Modell & Precision & Recall & Accuracy & F1 Score & Konsistenz \\
    \hline
    Opus 4.7   & 95,7 & 91,7 & 97,0 & 93,6 & [XX] \\
    Opus 4.8   & 90,9 & 83,3 & 94,0 & 87,0 & [XX] \\
    Opus 4.6   & 63,6 & 87,5 & 85,0 & 73,7 & [XX] \\
    Sonnet 4.6 & 62,5 & 83,3 & 84,0 & 71,4 & [XX] \\
    GPT 4.1    & 86,7 & 54,2 & 87,0 & 66,7 & [XX] \\
    GPT 5.5    & 47,5 & 79,2 & 74,0 & 59,4 & [XX] \\
    GPT 5.3    & 46,2 & 75,0 & 73,0 & 57,1 & [XX] \\
    GPT 5      & 44,8 & 54,2 & 73,0 & 49,1 & [XX] \\
    \hline
  \end{tabular}
\end{table}

Den höchsten F1 Score erreicht Opus 4.7 mit 93,6\,\%, den niedrigsten GPT 5 mit 49,1\,\%. Opus 4.7 und Opus 4.8 sind die einzigen Modelle, die sowohl bei der Precision als auch beim Recall über 80\,\% liegen. Opus 4.7 meldet dabei nur eine einzige unberechtigte Exception und übersieht zwei der 24 tatsächlichen Exceptions.

Bei den übrigen Modellen gehen Precision und Recall deutlich auseinander. Opus 4.6, Sonnet 4.6, GPT 5.5 und GPT 5.3 erreichen einen Recall zwischen 75,0\,\% und 87,5\,\%, ihre Precision liegt aber nur zwischen 46,2\,\% und 63,6\,\%. Diese Modelle melden also eher zu viele Exceptions. GPT 5.5 und GPT 5.3 kommen jeweils auf 21 \ac{FP} und damit auf fast so viele unberechtigte wie berechtigte Exceptions. Im Prüfalltag würde das für die Auditoren einen hohen Nachprüfaufwand bedeuten. GPT 4.1 zeigt das umgekehrte Muster. Mit einer Precision von 86,7\,\% sind die gemeldeten Exceptions meist berechtigt, mit einem Recall von 54,2\,\% wird aber fast die Hälfte der tatsächlichen Exceptions übersehen. GPT 5 schneidet bei beiden Kennzahlen schwach ab.

Die Accuracy ist in diesem Fall nur eingeschränkt aussagekräftig. Da nur 24 der 100 Rechnungen eine Exception enthalten, erreicht auch ein Modell wie GPT 4.1 trotz 11 übersehener Exceptions eine Accuracy von 87,0\,\%. Für den Vergleich der Modelle sind Precision, Recall und F1 Score daher aussagekräftiger.

\todo{Konsistenz kurz beschreiben: Welches Modell liefert bei wiederholten Durchläufen die stabilsten Ergebnisse?}

\section{Auswertung der Bewertungsmatrix}
\label{sec:auswertung_bewertungsmatrix}

Die Ergebnisse aus Abschnitt~\ref{sec:ergebnisse_stichproben} fließen nun in die Nutzwertanalyse ein. Die Kriterien und ihre Gewichtung wurden bereits in Kapitel~\ref{chap:evaluationsdesign} vor der Durchführung der Tests festgelegt und wurden danach nicht verändert. Für jedes Kriterium wurden die gemessenen Werte in Punkte auf einer Skala von [1 bis 5] überführt. \todo{Umrechnungsregel nennen} Multipliziert mit der jeweiligen Gewichtung ergibt sich der Nutzwert je Modell.

\begin{table}[htbp]
  \centering
  \caption{Bewertungsmatrix der Nutzwertanalyse}
  \label{tab:bewertungsmatrix}
  \begin{tabular}{lcccc}
    \hline
    Kriterium & Gewicht & [Modell A] & [Modell B] & [Modell C] \\
    \hline
    \hline
    Nutzwert & 100\,\% & [XX] & [XX] & [XX] \\
    \hline
  \end{tabular}
\end{table}

Um zu prüfen, wie stark das Ergebnis von der gewählten Gewichtung abhängt, wurde eine Sensitivitätsanalyse durchgeführt. Dabei wurden die Gewichte der Kriterien jeweils um [XX] Prozentpunkte verändert. \todo{Ergebnis, bleibt die Rangfolge stabil?}

\section{Vergleich der Modelle}
\label{sec:vergleich_modelle}

[Modell] erreicht mit [XX] Punkten den höchsten Nutzwert, gefolgt von [Modell] und [Modell]. \todo{Stärken und Schwächen je Modell in wenigen Sätzen gegenüberstellen, z.\,B. hohe Precision bei geringerer Konsistenz}
\section{Analyse der Ergebnisse}
\label{sec:analyse_ergebnisse}

\todo{Ergebnisse interpretieren und mit dem Stand der Forschung aus Kapitel~[XX] abgleichen. Mögliche Leitfragen: Decken sich die Ergebnisse mit bisherigen Studien? Welche Fehlerarten treten bei allen Modellen auf?}

Bei der Interpretation sind mehrere Einschränkungen zu beachten. Der Testdatensatz ist mit 100 Rechnungen eher klein, sodass bereits wenige abweichende Fälle die Kennzahlen deutlich verändern. Die Untersuchung fand ausschließlich innerhalb von Copilot Studio statt. Die beschriebenen Probleme beim Auslesen der Freigabescreenshots und bei PNG Dateien sind daher eher eine Grenze der Plattform als der Modelle selbst, beeinflussen aber die gemessenen Ergebnisse. Schließlich bilden die Ergebnisse den Stand der Modellversionen zum Zeitpunkt der Tests ab und können sich mit neuen Versionen ändern.

\section{Beantwortung der Forschungsfrage}
\label{sec:beantwortung_forschungsfrage}

Die Forschungsfrage dieser Arbeit lautet: \emph{Wie unterscheiden sich verschiedene LLMs hinsichtlich ihrer Eignung zur Unterstützung von Stichprobenanalysen im Audit-Bereich?}

Die Ergebnisse zeigen, dass sich die untersuchten Modelle vor allem in \todo{wichtigste Unterscheidungsmerkmale, z.\,B. Erkennungsrate, Konsistenz} unterscheiden. xyz eignet sich nach der Nutzwertanalyse am besten für den Einsatz im Internal Audit bei Freudenberg. \todo{Einschränken, in welcher Rolle, z.\,B. als Unterstützung bei der Vorprüfung, wobei die abschließende Bewertung weiterhin durch Auditoren erfolgt.}